\section{Problem Description}

\subsection*{Use case}

The course DAT255 covers the engineering side of deep learning: model selection, training,
evaluation, and deployment. A practical bottleneck for students taking the course is that
generic study tools, including general-purpose chatbots such as ChatGPT, are not grounded in the
DAT255 syllabus. They paraphrase well but cannot reliably point back to specific chapters of the
course material, and their answers drift between conventions used in different textbooks.

This project addresses that gap by building a teaching assistant grounded in the DAT255 reading
list. The assistant exposes four study activities through a single web interface:
explanations of concepts, multiple-choice quizzes, long-answer practice with feedback, and
flashcards. Each activity is conditioned on retrieved passages from the course book and
generated by a transformer language model.

\subsection*{Why deep learning}

The system has two functional requirements that together motivate a deep-learning solution.
First, the input is unstructured natural language, both for the user query and for the
underlying course material; classical retrieval methods such as TF--IDF can rank passages but
cannot rewrite them into a coherent explanation, generate plausible distractors for a quiz, or
produce a graded review of a free-text answer. Second, the output must follow four strict
formats (free-form text, JSON quiz, scored review, term/definition pair). Both components ---
dense retrieval over course chunks and conditional text generation in four task-specific
formats --- are well-studied applications of neural sequence models~\cite{vaswani2017attention,
lewis2020retrieval}.

\subsection*{System overview}

The solution combines two neural components: a retriever that finds relevant chunks of the
course book using sentence embeddings, and a generator that produces the answer in one of four
task formats. The two components are decoupled, which means the retriever can be re-indexed
when course material changes without retraining the generator. Figure~\ref{fig:system_overview}
shows the full pipeline.

\begin{figure}[h]
\centering
\begin{tikzpicture}[node distance=0.55cm and 0.7cm]

    % Top row: retrieval
    \node[data] (query)  {User query \\ \footnotesize(\texttt{<task>} + question)};
    \node[block, right=of query]  (embed)  {Query \\ embedding};
    \node[block, right=of embed]  (search) {HNSW \\ search};
    \node[data,  right=of search] (chunks) {Top-$k$ chunks};

    % Index above search
    \node[data, above=0.5cm of search] (index) {RAG index \\ \footnotesize(1\,960 chunks)};

    % Bottom row: generation, aligned with right half
    \node[data,  below=2.0cm of search] (prompt) {Prompt \\ \footnotesize(task + ctx + q)};
    \node[model, right=of prompt]       (gen)    {Decoder-only \\ Transformer};
    \node[data,  right=of gen]          (output) {Task output \\ \footnotesize(text/JSON)};

    % Top row edges
    \draw[arrow] (query)  -- (embed);
    \draw[arrow] (embed)  -- (search);
    \draw[arrow] (search) -- (chunks);
    \draw[arrow] (index)  -- (search);

    % Merge into prompt via a shared join point above the prompt
    \coordinate (join) at ($(prompt.north) + (0, 0.4)$);
    \draw[arrow] (query.south)  |- (join);
    \draw[arrow] (chunks.south) |- (join);
    \draw[arrow] (join) -- (prompt.north);

    % Bottom row edges
    \draw[arrow] (prompt) -- (gen);
    \draw[arrow] (gen)    -- (output);

\end{tikzpicture}
\caption{System overview. A user query is embedded and matched against the pre-built RAG index;
the top chunks are concatenated with the query and passed, prefixed by a task token, to the
generator.}
\label{fig:system_overview}
\end{figure}

\subsection*{Repository}

The full source code for this project is available at
\url{https://github.com/Hagestregen/dat255-teaching-assistant}.

\subsection*{Scope and contributions}

The mandatory course requirements are addressed as follows: the use case is identified above;
a decoder-only transformer is constructed and trained from randomly initialised weights
(Section~\ref{sec:model}); and four models are evaluated on held-out task-formatted data with
both intrinsic and extrinsic metrics (Section~\ref{sec:evaluation}). Beyond the mandatory
items, the project includes a comparison of four model strategies, retrieval-augmented
generation, hyperparameter logging through Weights \& Biases, and deployment as a Gradio web
application (Section~\ref{sec:deployment}).
